% ----------------------------------------------------------
% Metodologia
% ----------------------------------------------------------
\chapter{Gestão do Projeto}

\section{Organização da Equipe}

A equipe foi organizada de acordo com a identificação dos pontos positivos e das limitações de cada integrante, visando evitar entraves durante o desenvolvimento do projeto. Essa distribuição permitiu maior flexibilidade na execução das atividades e um aprofundamento mais amplo nos diferentes aspectos do tema proposto.

Essa abordagem teve como objetivo promover uma comunicação eficiente, facilitando a distribuição e o acompanhamento das tarefas, bem como a identificação ágil de problemas e a otimização dos recursos disponíveis ao longo do desenvolvimento.

\subsection{Responsabilidades e Papéis}

O Quadro~\ref{quad_equipe} apresenta a distribuição das funções atribuídas aos membros da equipe.

\begin{table}[H]
\caption{Composição da equipe e funções}
\label{quad_equipe}
\centering
\begin{tabular}{lccccc}
\hline
\textbf{Membro} &
\textbf{Gerente / PO} &
\textbf{Front End} &
\textbf{Back End} &
\textbf{Mobile} &
\textbf{Documentação} \\
\hline
Nathan    & X & - & - & - & X \\
Vinycius  & - & - & X & - & X \\
João      & - & - & - & X & X \\
Riquelmy  & - & X & - & - & X \\
\hline
\end{tabular}
\legend{Fonte: Elaborado pelos autores.}
\end{table}

O Gerente/Product Owner atua como o elo central da equipe, assegurando que as funcionalidades essenciais sejam implementadas de acordo com os objetivos do projeto. Entre suas atribuições estão a definição de prioridades do backlog, a organização das atividades de cada sprint e o acompanhamento do progresso do desenvolvimento.

O Desenvolvedor Front End é responsável pela construção da interface e da experiência visual da aplicação, utilizando tecnologias voltadas à apresentação e interação com o usuário, transformando os protótipos em telas funcionais e responsivas.

O Desenvolvedor Back End concentra-se na implementação da lógica de negócio, comunicação com serviços externos, gerenciamento de banco de dados e processamento das informações necessárias ao funcionamento do sistema.

O Desenvolvedor Mobile dedica-se à construção da aplicação para dispositivos móveis, garantindo compatibilidade, desempenho e usabilidade em diferentes tamanhos de tela e sistemas operacionais.

Por fim, os responsáveis pela documentação organizam e atualizam continuamente os registros do projeto, assegurando que decisões, processos e entregas permaneçam documentados para consulta futura.

\section{Metodologias de Gestão e Desenvolvimento}

A condução do projeto foi baseada em metodologias ágeis, que permitem maior flexibilidade e adaptação às mudanças durante o desenvolvimento. Entre as metodologias analisadas destacam-se o Scrum, o Feature Driven Development (FDD) e o PMBOK, cada uma apresentando características próprias para gerenciamento e acompanhamento de projetos.

No contexto do Linha Vital, optou-se pela utilização do Scrum devido à sua capacidade de organizar o trabalho em ciclos curtos e iterativos, denominados sprints, permitindo entregas contínuas e incrementais. Essa abordagem favorece a comunicação entre os membros da equipe, a rápida identificação de problemas e a adaptação constante às necessidades do projeto \cite{scrumguide2020}.

\subsection{Scrum}

O Scrum é uma metodologia ágil que estrutura o desenvolvimento de projetos em ciclos de trabalho denominados sprints, caracterizados por duração previamente definida e objetivos específicos. Ao término de cada sprint é entregue um incremento funcional do produto, permitindo validações frequentes e correções ao longo do processo de desenvolvimento.

Essa metodologia promove transparência, inspeção contínua e adaptação, princípios fundamentais para projetos que exigem flexibilidade e evolução constante \cite{scrumguide2020}.

\subsubsection{Sprints}

No projeto Linha Vital, os sprints foram planejados para contemplar atividades como levantamento de requisitos, modelagem do sistema, implementação das funcionalidades, realização de testes e elaboração da documentação.

Ao final de cada ciclo, a equipe realizava reuniões de revisão e retrospectiva, avaliando os resultados obtidos e identificando oportunidades de melhoria para os próximos períodos de desenvolvimento. Essa prática garantiu maior colaboração entre os integrantes e contribuiu para a evolução contínua do projeto.

\section{Repositório da Aplicação}

O repositório da aplicação foi definido como o ambiente central para armazenamento e versionamento do código-fonte, da documentação e dos demais artefatos produzidos durante o desenvolvimento.

Para essa finalidade, foi adotado o GitHub, plataforma amplamente utilizada no desenvolvimento de software por disponibilizar recursos de controle de versão, colaboração entre equipes e rastreabilidade das alterações realizadas \cite{progit2014, githubdocs2026}.

\subsection{Definição do Repositório da Aplicação}

No contexto do Linha Vital, o repositório foi estruturado de forma a organizar os diferentes componentes do sistema, separando os módulos de desenvolvimento, documentação e demais recursos do projeto.

Essa organização facilitou a manutenção do código, o gerenciamento das versões e o trabalho colaborativo entre os integrantes da equipe. Além disso, a utilização do GitHub permitiu manter um histórico completo das alterações realizadas ao longo do desenvolvimento, contribuindo para a transparência e confiabilidade do projeto.

\subsection{Link do Repositório e Especificações para Acesso}

O repositório encontra-se hospedado na plataforma GitHub, que oferece mecanismos de controle de versão, rastreabilidade das alterações e colaboração entre desenvolvedores \cite{githubdocs2026}.

O acesso ao ambiente foi configurado de forma restrita aos membros da equipe e aos orientadores do projeto, mediante autenticação individual e controle de permissões. Essa configuração garante maior segurança e permite o acompanhamento adequado das contribuições realizadas durante o desenvolvimento.

Link do repositório:

\begin{center}
\url{https://github.com/Vinycius-Silva/Projeto_Conclusao_Linha_Vital}
\end{center}

Especificações de acesso: credenciais individuais fornecidas pela equipe de desenvolvimento, com permissões diferenciadas para colaboradores e administradores.